\section{Introduction}
\label{sec:introduction}

Polynomial arithmetic is a central computational component of modern
zero-knowledge proof systems. Provers manipulate large polynomials during
arithmetization, constraint evaluation, low-degree testing, polynomial
commitment, quotient construction, and opening procedures. Transparent
proof systems based on Reed--Solomon proximity testing, including
FRI-based constructions and STARKs, make polynomial evaluation and encoding
particularly visible components of the prover computation
\cite{stark2018}.

The representation chosen for a polynomial determines the cost of many of
these operations. A polynomial may be represented by its monomial
coefficients, by its values on a structured evaluation domain, or by
coefficients in another basis adapted to a particular algebraic operation.
These representations describe the same polynomial space but expose
different arithmetic structure and different opportunities for
parallelism.

This work isolates one basic operation:
\begin{quote}
Given a degree-\(<N\) univariate polynomial \(P\) and an arbitrary point
\(x\) in the base field, compute \(P(x)\) from the chosen representation.
\end{quote}

We deliberately study this primitive independently of a complete proof
system. This allows the effect of the representation itself to be measured
under a common field implementation, machine, input size, and benchmarking
methodology.

Let
\[
N=2^m.
\]
In addition to the monomial representation and a roots-of-unity
Lagrange/evaluation representation, we study the following family of
univariate polynomials:
\[
K_y(X)
=
\prod_{i=0}^{m-1}
\left(X^{2^i}+y_i\right),
\qquad
y=(y_0,\ldots,y_{m-1})\in\{0,1\}^m.
\]
We refer to this family as the \emph{Boolean-kernel basis}. A polynomial
represented in this basis is written
\[
P(X)
=
\sum_{y\in\{0,1\}^m}
c_y K_y(X).
\]

The basis admits a simple binary reduction rule. Fix an evaluation point
\(x\), and define
\[
t_i=x^{2^i}.
\]
At level \(i\), two coefficients \(a\) and \(b\) whose indices differ only
in bit \(i\) are replaced by
\[
F_{t_i}(a,b)
=
t_i(a+b)+b.
\]
Applying this rule recursively reduces \(N\) coefficients to the single
value \(P(x)\). The full tree contains exactly \(N-1\) folds, and all
folds at the same level are independent.

This structure makes the evaluation algorithm both linear in the number of
input coefficients and naturally parallel. Since an arbitrary coefficient
vector must in general be inspected in full, an \(O(N)\) algorithm already
has optimal asymptotic input complexity for this task. The relevant
question is therefore whether the basis gives a practically favorable
constant factor and memory-access pattern.

We answer this experimentally using a Rust implementation over the
Goldilocks field
\[
\mathbb{F}_p,
\qquad
p=2^{64}-2^{32}+1.
\]
The optimized implementation combines coarse-grained subtree parallelism,
a specialized Goldilocks reduction, a fused kernel-fold operation,
bounds-check-free inner loops, direct first-level reduction, and
worker-local reusable scratch storage.

For
\[
N=2^{20}=1{,}048{,}576,
\]
the final same-run cross-basis benchmark on the tested Apple M1 machine
gives a central estimate of
\[
\boxed{1.0560\,\mathrm{ms}}
\]
for the optimized Boolean-kernel evaluator and
\[
1.4116\,\mathrm{ms}
\]
for the blocked parallel monomial evaluator.  Thus the kernel evaluator is
\[
\boxed{1.337\times}
\]
faster in this controlled comparison, corresponding to approximately
\[
\boxed{25.2\%}
\]
less wall-clock evaluation time.  During kernel-only optimization, a
separate run reached \(892.57\,\mu\mathrm{s}\); we report that value as a
best observed implementation result but do not mix it with competitor
measurements from another run when computing cross-basis speedups.

This comparison is intentionally an operation-level result. It does
\emph{not} imply a \(1.337\times\) speedup for a complete STARK, FRI prover,
or polynomial commitment scheme. A complete protocol also contains
encoding, low-degree extension, hashing, commitment, transcript, query,
and other polynomial operations, and a change of representation may create
conversion costs elsewhere. The result instead establishes that
representation alone can materially change the cost of arbitrary-point
evaluation and motivates studying whether this basis can be maintained
across a larger fraction of a prover pipeline.

Related work on transparent proof systems already demonstrates the
importance of linear- and quasi-linear-time polynomial machinery
\cite{stark2018}. Other polynomial commitment constructions have also
highlighted the cost of moving between polynomial representations; for
example, Zeromorph explicitly addresses interactions between multilinear
polynomials and univariate commitment machinery \cite{zeromorph2023}.
Our scope is different: we experimentally isolate a univariate basis and
measure its effect on a concrete polynomial operation.

\paragraph{Contributions.}
The contributions of this work are:
\begin{enumerate}
    \item We give a direct binary-fold evaluation algorithm for the
    univariate Boolean-kernel basis and derive its exact arithmetic cost.

    \item We provide an optimized Rust implementation over the Goldilocks
    field with correctness tests and reproducible benchmarks.

    \item We compare single-point evaluation against conventional
    univariate representations under a common implementation methodology,
    including a parallel monomial baseline rather than an artificially
    serialized competitor.

    \item We provide an optimization ablation covering arithmetic reduction,
    memory organization, bounds checks, scratch-buffer reuse, subtree size,
    and unsuccessful tree-fusion strategies.

    \item We distinguish primitive-level speedups from end-to-end
    proof-system claims and identify the representation-conversion and
    integration questions that must be resolved before transferring the
    result to a complete prover.
\end{enumerate}

The remainder of the paper first defines the polynomial representations
formally, then derives the Boolean-kernel evaluation algorithm, describes
the implementation and benchmark methodology, reports the experimental
results, and discusses implications for proof-system design.
